\documentclass[14pt,a4paper]{extarticle}
\usepackage[T2A]{fontenc}
\usepackage[utf8]{inputenc}
\usepackage[english,russian]{babel}
\usepackage{lmodern}
\usepackage{geometry}
\usepackage{array}
\usepackage{tabularx}
\usepackage{longtable}
\usepackage{ltablex}
\usepackage{booktabs}
\usepackage{hyperref}
\usepackage{indentfirst}
\usepackage{microtype}
\usepackage{graphicx}
\usepackage{setspace}
\usepackage{listings}
\usepackage{graphicx}
\usepackage{caption}
\usepackage{float}
\usepackage{tocloft}
\usepackage{titlesec}
\usepackage{newtxtext,newtxmath}

\titleformat{\section}        % форматируем \section
  {\normalfont\bfseries}      % стиль: обычный + жирный
  {\thesection.}              % номер секции с точкой
  {2mm}                       % расстояние между номером и текстом
  {}                           % перед заголовком можно добавить что-то
\titlespacing*{\section}{1.5cm}{*3}{*1}  % {отступ слева}{отступ сверху}{отступ снизу}

\captionsetup{justification=centering}
\renewcommand{\ttdefault}{cmtt}
\keepXColumns
\setlength{\parindent}{3em}
\newgeometry{left=2.5cm,right=2.5cm,top=3cm,bottom=2.5cm}

\renewcommand{\contentsname}{\hspace*{4em}{\smallСодержание}}
\setlength{\cftbeforesecskip}{0.5em}
\renewcommand{\cftsecleader}{\cftdotfill{\cftdotsep}}
\renewcommand{\cftsecaftersnum}{.\hspace{2mm}}

\title{Отчёт по "Программно-аппаратные средства интеллектуальных систем"}
\author{Выполнил студент группы М-РИТ-25-1, Станиславчук Сергей}

\begin{document}

\setstretch{1.5}
\noindent
\begin{minipage}[t]{0.01\textwidth}
    \includegraphics{~/programmer/++Programmer/5_1/LSTU.jpg}
\end{minipage}%
\begin{minipage}[c]{1.1\linewidth}
    \centering
    {\small
    МИНИСТЕРСТВО НАУКИ И ВЫСШЕГО ОБРАЗОВАНИЯ РФ \\[1em]
    ФЕДЕРАЛЬНОЕ ГОСУДАРСТВЕННОЕ БЮДЖЕТНОЕ \\
    ОБРАЗОВАТЕЛЬНОЕ УЧРЕЖДЕНИЕ ВЫСШЕГО ОБРАЗОВАНИЯ \\
    «ЛИПЕЦКИЙ ГОСУДАРСТВЕННЫЙ ТЕХНИЧЕСКИЙ УНИВЕРСИТЕТ»
    }
\end{minipage}%

\vspace{3em}

\noindent
\hspace{1em}%
\begin{minipage}[t]{0.8\textwidth}
\textbf{Институт} \hspace{2em} компьютерных наук \\
\textbf{Кафедра} \hspace{2em} автоматизированных систем управления
\end{minipage}


\vspace{6em}

\begin{center}
\textbf{Практическая работа №2} \\
по дисциплине \\ <<Планирование карьеры>>
\end{center}

\vspace{6em}

\noindent
\begin{tabular*}{\textwidth}{@{\extracolsep{\fill}} l c r }
    Студент М-РИТ-25-1 & \makebox[4cm]{\hrulefill} & Станиславчук С.М. \\[2em]
Руководитель & \\
Доцент & \makebox[4cm]{\hrulefill} & Триханова И.Е. \\
\end{tabular*}

\vfill
\begin{center}
    Липецк, 2026
\end{center}

\thispagestyle{empty}

\newpage
\tableofcontents
\newpage

\section{Вопрос 27. Какие навыки сегодня ценятся выше: hard skills или soft skills?}

С точки зрения работодателя, оценка значимости hard skills и soft skills зависит от контекста рынка труда и изменений в способах получения знаний. Если рассматривать ситуацию ретроспективно, то около 30 лет назад приоритет в большей степени отдавался именно hard skills. Это было обусловлено ограниченностью источников информации: профессиональные знания приобретались преимущественно либо в образовательных учреждениях, либо непосредственно на рабочем месте. В обоих случаях процесс обучения неизбежно включал активное взаимодействие с другими людьми — преподавателями, наставниками, коллегами. Таким образом, развитие soft skills происходило как естественный побочный эффект профессионального становления, даже если оно не являлось осознанной целью.

В современных условиях ситуация изменилась принципиально. Доступ к информации стал практически неограниченным благодаря интернету. В частности, в таких областях, как компьютерные науки, значительная часть знаний может быть получена без прямого взаимодействия с другими людьми: через документацию, онлайн-курсы, форумы и специализированные ресурсы. Это позволяет эффективно развивать hard skills в изолированном формате, однако одновременно снижает необходимость в регулярной коммуникации в процессе обучения. В результате специалист может обладать высоким уровнем технической подготовки, но испытывать дефицит коммуникативных и социальных навыков.

При этом рабочая среда по-прежнему требует взаимодействия: участие в командной разработке, обсуждение решений, согласование задач и работа с заинтересованными сторонами остаются неотъемлемыми элементами профессиональной деятельности. Вследствие этого работодатели начинают уделять повышенное внимание soft skills, поскольку именно они оказываются недостаточно развитыми у части кандидатов. Создается впечатление, что эти навыки ценятся выше, однако это скорее отражение дефицита, чем реального смещения приоритетов.

Таким образом, с позиции работодателя hard skills и soft skills следует рассматривать как взаимодополняющие компоненты профессиональной компетентности. Недостаток любой из этих составляющих снижает ценность специалиста на рынке труда. Несмотря на временный акцент на развитии soft skills, обусловленный изменениями в образовательной и информационной среде, в действительности обе группы навыков имеют сопоставимую значимость. Специалист, не обладающий достаточным уровнем как технических, так и коммуникативных компетенций, с высокой вероятностью не будет соответствовать требованиям современной профессиональной среды.

\newpage
\section{Вопрос 37. Можно ли быть продуктивным без чёткого расписания?}

С точки зрения разработки программного обеспечения и проектной работы в целом, продуктивность без жёсткого расписания возможна. Более того, в ряде случаев именно отказ от строгой регламентации рабочего времени позволяет достичь более высоких результатов.

Прежде всего, следует разграничить понятие расписания как инструмента координации и как средства индивидуальной организации труда. Расписание исторически формировалось как механизм синхронизации деятельности группы: оно необходимо в образовательных учреждениях и других системах, где требуется одновременное участие большого количества людей в общем процессе. В профессиональной среде, особенно в проектной разработке, задачи различных специалистов существенно различаются, а их выполнение не всегда требует строгой временной привязки.

Профессия программиста в её современном виде сформировалась в конце XX века, с
развитием персональных компьютеров и сетевых технологий. В этой сфере уже
длительное время применяется практика гибкого графика. Как правило, он включает
фиксированные интервалы для синхронизации команды (например, совещания или
статус-митинги), тогда как остальное рабочее время остаётся в распоряжении
сотрудника. Такой подход позволяет разработчику самостоятельно определять
периоды наибольшей концентрации и распределять нагрузку в соответствии с
индивидуальной продуктивностью.

При этом результативность труда в значительной степени определяется не
количеством отработанных часов, а степенью вовлечённости в задачу. Конечно же
можно решить все необходимые задачи в начале рабочей недели и решить - сказать
своему боссу о выполненной работе, чтобы он добавил новой (за вознаграждение,
разумеется), либо остановить свой рабочий процесс и дождаться добавления задач
по расписанию. Это, на мой взгляд, является отличным балансом между работой и
отдыхом, что снижает риск эмоционального выгорания.

В противоположность этому, в более консервативных трудовых культурах, например
в Японии, по-прежнему распространены жёсткие графики и высокая степень
формализации рабочего времени. Это нередко приводит к переработкам и снижению
общей эффективности труда из-за накопления усталости. Таким образом, отсутствие
строгого расписания при наличии минимальной координации внутри команды может не
только не снижать продуктивность, но и способствовать её росту за счёт
повышения автономии и мотивации специалиста.


\newpage
\section{Вопрос 43. Какие ошибки чаще всего допускаются при постановке целей?}
Ограничение по времени (сроки). Часто при постановке целей закладываются излишне оптимистичные сроки выполнения задач. Возникает стремление минимизировать время реализации, что приводит к недооценке сложности работы. В результате графики сдвигаются, а отдельные элементы задачи остаются недоработанными или требуют последующей доработки. Это снижает общее качество результата и нарушает планирование последующих этапов проекта.

Недостаточная квалификация специалистов. Цели могут формулироваться без учёта текущего уровня компетенций исполнителей. В таких условиях задача оказывается чрезмерно сложной для команды, что приводит к задержкам, ошибкам и снижению мотивации сотрудников. Корректная постановка целей требует соотнесения сложности задачи с реальными возможностями команды.

Отсутствие чёткой формулировки цели. Нечёткие, размытые или двусмысленные цели затрудняют их интерпретацию исполнителями. Это приводит к тому, что разные участники проекта по-разному понимают ожидаемый результат, что увеличивает риск несоответствия итогового продукта изначальным ожиданиям.

Игнорирование измеримости результата. Если цель не содержит критериев оценки, становится невозможно объективно определить, достигнута она или нет. Отсутствие метрик затрудняет контроль выполнения и снижает управляемость процесса.
\end{document}
